\begin{abstract}
本文给出一个以 Zeckendorf 截断折叠为核心的“信息守恒完备系统”实现。空间状态是一个长度为 \(m\) 的比特串 \(y\in Y_m\)（无相邻 \(11\)），时间纤维（uplift 坐标）是窗口外尾部比特串 \(t\in T_m\)。观察者的动力学只包含两类互逆操作：\textbf{展开}与\textbf{折叠}。展开每推进一步都会在 \(\Tail^{-1}\) 诱导的可行候选中做出一次\textbf{可审计分支选择}，从而产生新的时间纤维/微观解释候选；同时，由于折叠核 \(\Fold\) 是多对一，空间侧的纤维大小 \(s_m(y)=|\Fold^{-1}(y)|\) 提供可采掘的简并度，我们把它量化为潜在能量 \(\Energy_{\mathrm{pot}}(y)=\log_2 s_m(y)\)。观察者并行维护候选分支集合 \(\mathcal B\) 的规模受一个\textbf{资源上限} \(W\) 约束（beam width 上限，定义 \ref{def:energy_beam}）；折叠则以展开的逆操作逐步湮灭对应候选并回到先前状态。为保证互逆性，观察者携带一条 trace \(\mathsf{tr}\in\{0,1\}^\ast\) 作为可审计的路径标签。我们实现了可复现实验：在 \(m=6\) 下枚举并统计纤维大小 \(2/3/4\) 的分布，并在固定观测 \(y_t=0^m\) 下展示展开深度与资源上限约束下的分支行为。
\end{abstract}

\paragraph{贡献点（最小主张清单）}
\begin{itemize}
  \item 给出以 Zeckendorf 截断折叠 \(\Fold\)（\(2^m\to F_{m+2}\)）为\textbf{空间压缩核}、以窗口外尾部为\textbf{uplift 时间纤维}的完备系统定义，并把本体固定为静态约束 \(\Clo\)。
  \item 用抽象时间算子 \(\Tail\) 及其逆候选 \(\Tail^{-1}\) 给出明确的 uplift 规则；并把观察者并行维护候选集合的限制刻画为资源上限 \(W\)，从而把“多分支与主观单线”落实为可审计的截断（commit）规则。
  \item 把折叠核的纤维大小 \(s_m(y)=|\Fold^{-1}(y)|\) 作为空间状态的结构属性，并定义潜在能量 \(\Energy_{\mathrm{pot}}(y)=\log_2 s_m(y)\)；在基例 \(m=6\) 的 \(64\to 21\) 核中得到三值谱 \(s_6(y)\in\{2,3,4\}\) 与可复现的分布直方图。
  \item 以资源上限 \(W\) 给出 beam width，并以此把“多世界/分裂/主观选择”翻译为可审计的资源约束结论。
  \item 提供一键复现实验流水线，生成 `sections/generated/` 可直接引用的图表片段与稳定导出物 `artifacts/export/`。
\end{itemize}

\tableofcontents
\newpage

